\chapter{Multi-Label Classification and Class-Balanced Sampling (Phase 1.2)}
\label{chap:p1-2}

\section{Goal}

The pilot assigned each image exactly one class by \emph{priority collapse}: a row that violated more
than one rule was labeled only with its highest-priority violation
(ppe\_violation $>$ fall\_hazard $>$ struck\_by\_risk), and every lower-priority violation on the same
image was discarded. This is a labeling convention, not a property of the dataset, and it has a
concrete, damaging consequence for the scarcest and most safety-critical class. An image showing a
worker with no hard hat \emph{and} standing in an excavator's swing radius is a real struck-by hazard,
but priority collapse files it as a PPE violation and the struck-by label is lost. The pilot's
struck\_by\_risk class was therefore powered at only $n=13$ in the entire test split --- too small for
any stable estimate.

The scale-up plan (item 3.2) diagnoses this precisely: the struck\_by\_risk scarcity is ``a labeling
artifact, not a dataset property,'' and predicts that multi-label classification recovers the hidden
examples. The goal of this phase is to fix the labeling policy so that an image counts toward --- and
is tested against --- every rule it actually violates, and to verify the predicted recovery.

\section{What Was Done}

\subsection{Exposing the full violated-class set}

The information priority collapse discards was already being computed and thrown away. A new function,
\texttt{data.image\_classes(row)}, returns the \emph{complete} list of violated classes for a row
(highest-priority first), where the existing \texttt{classify\_image} returns only the single collapsed
primary class. A clean row returns \texttt{["compliant"]}; a rule\_1+rule\_4 row returns
\texttt{["ppe\_violation", "struck\_by\_risk"]}; a rule\_2+rule\_3 row returns \texttt{["fall\_hazard"]}
once, since both those rules map to the same class.

\subsection{A labeling-aware sampler}

\texttt{select\_balanced\_sample} gained a \texttt{labeling} parameter:

\begin{itemize}
  \item \texttt{"priority"} (default): the frozen pilot sampler, untouched in behavior --- one class
  and one rule per image. Its only addition is that it now also populates the new
  \texttt{assigned\_rule\_ids} field with the single-element list \texttt{[assigned\_rule\_id]}, so
  that both labeling modes present the same interface downstream.
  \item \texttt{"multilabel"}: an image is placed into the candidate reservoir of \emph{every} class it
  violates, so a rule\_1+rule\_4 image sits in both the ppe\_violation and struck\_by\_risk pools and
  can be drawn to satisfy either class's coverage floor. After per-class selection, images drawn via
  more than one class are deduplicated by \texttt{image\_id} and their violated rules unioned, so each
  unique image appears once carrying \emph{all} the rules it is tested against.
\end{itemize}

The \texttt{PilotSample} dataclass gained an \texttt{assigned\_rule\_ids: list[str]} field to hold that
per-image rule set. Determinism is preserved exactly: multi-label selection uses the same seeded
shuffle in the same class order as the priority sampler, and deduplication preserves first-encounter
order.

\subsection{Manifest and configuration}

\texttt{config.LABELING} holds the default policy (\texttt{"priority"}). Script
\texttt{02\_select\_samples.py} gained a \texttt{-{}-labeling} flag. Under priority it writes the
original five-column \texttt{pilot\_samples.csv} unchanged; under multi-label it writes an expanded
six-column, mode-suffixed manifest (\texttt{pilot\_samples\_multilabel.csv}) with an
\texttt{assigned\_rule\_ids} column, so the two policies never share an output file.

\section{Results}

Scanning the full 3{,}004-image test split (reading only the four rule columns from the local parquet,
no image decoding), the predicted recovery appears exactly as the plan described it:

\begin{center}
\begin{tabular}{p{4.4cm}p{3.6cm}p{3.6cm}}
\toprule
\textbf{class} & \textbf{priority (primary)} & \textbf{multi-label (coverage)} \\
\midrule
ppe\_violation   & 324 & 324 \\
fall\_hazard     & 75  & 86 \\
struck\_by\_risk & \textbf{13}  & \textbf{24} \\
compliant        & 2{,}592 & 2{,}592 \\
\bottomrule
\end{tabular}
\captionof{table}{Per-class counts on the full test split. Priority-collapse counts each image once,
under its highest-priority class; multi-label counts an image toward every class it violates.}
\label{tab:p12-recovery}
\end{center}

The struck\_by\_risk pool nearly doubles, from 13 to 24 --- a recovery of exactly 11 images. Those 11
are precisely the images that violate rule\_4 while \emph{also} carrying a higher-priority violation
that priority-collapse had used to file them elsewhere; the count was verified directly, not inferred:
11 struck-by images have a non-struck-by primary class under the old policy. This reproduces to the
image the pilot report's own statement that ``11 of 24 real struck\_by\_risk test-split images also had
a higher-priority violation.'' fall\_hazard likewise recovers 11 images (75 $\rightarrow$ 86) hidden
behind PPE violations. ppe\_violation and compliant are unchanged, as expected: PPE is the top
priority (nothing outranks it to hide it) and compliant images have no violations to redistribute.

The full unit suite is green at 76 tests (68 after Phase~1.1, plus 7 new multi-label tests written
test-first, plus one already counted). The seven new tests pin the full-set classification, the
fall\_hazard de-duplication, the per-image multi-rule assignment, the struck-by recovery, and the
guarantee that a shared image is never duplicated across class reservoirs.

\section{A Scoping Decision, Stated Plainly}

This phase deliberately fixes the \emph{data layer} --- classification, sampling, and per-image rule
assignment --- and stops there. It does \emph{not} rewire the downstream metric scripts (baseline
inference through completeness) to iterate over \texttt{(image, rule)} test units rather than one row
per image. That rewiring is real work with real blast radius: every metric script currently keys on
\texttt{image\_id} with a single assigned rule, and \texttt{baseline\_predictions.csv} is one row per
image.

The decision to defer it is not a shortcut; it follows the plan's own structure. The scale-up plan
couples the multi-label fix directly with the Phase~5 scale run (``scale from 163 to the full
3{,}004-image test split \ldots pair with the multi-label fix so struck\_by\_risk is powered''). Doing
the end-to-end rewiring now would fork the frozen 163-sample set in the middle of the Phase~1
validation sequence, where every other fix is being checked \emph{against} that same frozen set.
Building the recovery mechanism now, and consuming it at scale later, keeps each Phase~1 fix
independently verifiable against a stable baseline. The consumption path is a Phase~5 deliverable, and
is called out as such here so the boundary is explicit rather than implied.

\section{Standard vs. Non-Standard Choices}

\textbf{Standard.} Multi-label treatment of images that genuinely exhibit multiple hazards is the
correct and conventional framing for a multi-hazard detection problem; single-label collapse was only
ever a simplification the pilot adopted to get a balanced four-class subset quickly.

\textbf{Non-standard, and worth noting.} The recovery was validated by reading only the rule columns
from the parquet and counting, rather than by running the full sampler over decoded images --- a
deliberately cheap check that isolates the one claim being tested (does multi-label surface the hidden
struck-by images?) from every slower, unrelated concern. And, consistent with Phase~1.1, the new
policy is gated behind a default-valued flag that writes to a separate manifest file, so the frozen
pilot manifest and every number derived from it remain reproducible while the multi-label manifest
exists alongside it.

\section{Implications}

struck\_by\_risk is no longer inherited at $n=13$; on the test split alone it is powered at $n=24$, and
across the train split (available for supplementing minority classes) the plan's target of roughly 70
usable examples is now reachable. More importantly, multi-label data is the precondition for the
multi-hazard completeness tests the proposal promised (Phase~2.5): an image that violates two rules can
now be labeled, sampled, and --- once the Phase~5 consumption path lands --- tested for whether the
explanation captures \emph{both} hazards. The immediate carry-forward is that boundary: the sampler and
manifest speak multi-label now, and the metric scripts will learn to consume it during the scale run.
